% Regulatory Analysis subsection — standalone LaTeX fragment
% Intended to replace or expand \subsection{Regulatory Implications} in universal_impossibility.tex
% FOCUSED on SR 11-7, with brief notes on EU AI Act and NIST AI RMF

\subsection{Regulatory Analysis: SR~11-7}
\label{sec:regulatory-analysis}

OCC SR~11-7 \citep{occ2011sr117} establishes three pillars for model
validation in banking that are analogous to the trilemma.
SR~11-7's discussion of \emph{conceptual soundness} requires that ``the
model's theoretical construct, \ldots\ assumptions, and limitations should
be assessed''---this parallels \emph{faithfulness}: the explanation must
reflect the model's actual structure and assumptions.
SR~11-7's \emph{outcomes analysis requirements} call for ``comparing the
model's outputs with the corresponding realized outcomes'' across ``model
versions''---this parallels \emph{stability}: the explanation must be
consistent across model versions that produce equivalent outcomes.
SR~11-7's \emph{effective challenge framework} requires ``critical analysis
by objective, informed parties'' yielding ``specific findings''---this
parallels \emph{decisiveness}: the challenge must produce specific,
actionable findings, not hedged non-answers.

These three requirements, as stated, cannot be simultaneously satisfied for
any banking model with the Rashomon property---i.e., any model where
multiple parameter configurations achieve equivalent predictive
performance. For models trained with bagging, boosting, or stochastic
gradient descent, the Rashomon property is generic
(Section~\ref{sec:resolution}). Theorem~\ref{thm:universal} applies
directly.

Compliant model documentation under the trilemma should: (a)~state which
property is prioritized for each model use case; (b)~use ensemble
explanations (\DASH{}) for the stable $+$ faithful tradeoff, documenting
where decisiveness is sacrificed; (c)~report explanation uncertainty bounds
(confidence intervals on feature importance across the Rashomon set)
alongside point estimates; and (d)~document the Rashomon set size and the
resulting explanation variance as part of the model risk assessment.
SR~11-7's effective challenge framework is then directed at modeling choices---why
this Rashomon set? why this loss threshold?---rather than at explanation
artifacts that are mathematically underdetermined.

The same analysis applies to EU AI Act Article~13 \citep{euaiact2024}
(transparency) and NIST AI RMF MAP~1.5 \citep{nistai2023}
(interpretability). In each case, the regulatory language implicitly
assumes all three properties are jointly achievable; our theorem shows they
are not. Regulations should be reformulated as explicit tradeoffs: specify
which property is required, which may be relaxed, and what resolution
method is acceptable.
